\section{A native source certificate for the chamber graph}
\label{sec:chambers}

The cographic persistence proposal in repository report~92~\cite{persistence} assumes an
authenticated chamber graph. Its integration contract explicitly leaves the
native source exporter open. We now implement this interface for the
components of one supplied embedded normal surface. The new operator is
\texttt{normal\_chamber\_graph}; independent replay is
\texttt{verify\_normal\_chamber\_graph}. Both operate on the maintained finite,
connected, orientable source triangulations with one torus boundary.

This section establishes component identities before using them as graph
edges. A list of compatible normal coordinate vectors does not by itself
specify labelled disjoint representatives. Our input instead consists of one
normal vector $x$ satisfying the source matching and quadrilateral constraints,
together with one proposed representative disc
\[
 d_i=(t_i,c_i,r_i),\qquad 0\le r_i<x_{t_i,c_i},\qquad 0\le i<k,
\]
per component. The connected components of this single embedded surface
already have a common geometric realization. Exact orbit certificates verify
that the proposed labels are exhaustive and distinct.

\subsection{Authenticating the component labels}

Let $X$ be the maintained normal-arc interval system. Its points are normal
intersection points on globally oriented source edges; its orbits are the
connected components of the normal surface. For each marked disc choose one
of its normal vertices $p_i\in X$. The source chart authenticates that this
point belongs to $d_i$.

\begin{lemma}[Exhaustive labels from a scalar census]
\label{lem:component-marker-census}
Assign the weight $2^i$ to $p_i$, adding weights when marked points coincide.
The exact weighted orbit histogram equals
\[
 \bigl\{(2^i,1):0\le i<k\bigr\}
\]
if and only if the marked discs lie in distinct surface components and every
component contains a marked disc.
\end{lemma}
\begin{proof}
Each marker is assigned once. An orbit has weight equal to the sum of the
distinct marker powers belonging to it, so the sum has no binary carries.
An unmarked component contributes an orbit of weight zero. A component
containing two or more labels contributes a weight with at least two nonzero
bits. The stated histogram therefore forces one marker per component and no
unmarked component. Conversely, an exhaustive distinct labelling gives
exactly those $k$ singleton-bit weights, each with multiplicity one.
\end{proof}

Two different discs on one component are consequently rejected as a proposed
two-component labelling. A component with no representative cannot disappear
from the output. The producer performs this scalar weighted query first and
returns \texttt{INCONCLUSIVE} without a graph if the required histogram fails.
When no representative list is supplied, it tries the first present normal
disc. This deliberate default handles a connected nonempty surface; it does
not assert that an arbitrary supplied vector is connected.

\begin{lemma}[The two-sidedness gate]
\label{lem:component-two-sided}
Suppose Lemma~\ref{lem:component-marker-census} authenticates $k$ components.
All components are two-sided if and only if the normal surface with vector
$2x$ has $2k$ components.
\end{lemma}
\begin{proof}
The surface with doubled coordinates is the horizontal boundary of the
normal interval bundle. A connected two-sided component contributes two
horizontal components, whereas a connected one-sided component contributes
one connected double. If their respective counts are $a$ and $b$, the two
component counts are $a+b=k$ and $2a+b$. The latter equals $2k$ precisely
when $b=0$.
\end{proof}

The implementation directly checks the doubled orbit count. It does not
divide by a common coordinate factor. An odd multiplicity of a one-sided
surface can contain a one-sided component; its even double can instead be
connected and two-sided. Both behaviours are covered by native fixtures.
For every accepted connected component, choosing a side at the representative
disc determines a global coorientation. Choosing a different disc can reverse
that coorientation, which simply reverses the corresponding graph edge.

\subsection{Locating the sides in the actual cut}

The cut-chamber interval system identifies local normal chambers only through
original face regions. It never crosses the supplied surface. Its orbits
are therefore the connected components of the manifold cut along all the
surface components.

The marker addresses use the existing source rank conventions. In a
tetrahedron containing $q$ quadrilaterals, the central chamber chain has
addresses $c,c+1,\ldots,c+q$. Quadrilateral number $r$ lies between $c+r$
and $c+r+1$. A triangle stack at a vertex has consecutive arm chambers;
each triangle lies between consecutive arm chambers, except the last one,
which meets the appropriate end of the central chain. These formulas use
binary integer ranks and do not enumerate the intervening chambers.

\paragraph{The affine addresses used by the exporter.}
For precision, write the triangle counts in one tetrahedron as
$t_0,\ldots,t_3$. If quadrilaterals are present, their unique allowed
column is $4+j$, with $j\in\{0,1,2\}$, and their number is $q$.
When $q=0$, set $j=0$ purely for the formulas below.
In the maintained convention, this type separates the vertex pair
$\{0,j+1\}$ from its complementary pair. Let $C$ be the first central
chain address. The first arm address at vertex $v$ is
\[
 A_v=C+q+1+\sum_{u<v}t_u.
\]
The next tetrahedron starts after the chain and all four arms.
The adjacent chamber atoms of triangle $(v,r)$ are
\[
 \left(A_v+r,\
 \begin{cases}
 A_v+r+1,&r+1<t_v,\\
 C,&r+1=t_v\text{ and }(q=0\text{ or }v\in\{0,j+1\}),\\
 C+q,&r+1=t_v,\ q>0,\ v\notin\{0,j+1\}.
 \end{cases}\right).
\]
For quadrilateral $(4+j,r)$ they are simply $(C+r,C+r+1)$.
These pairs are ordered choices of local negative and positive side.

The surface-component marker also has an explicit affine address.
For a triangle at vertex $v$, choose the directed local edge from $v$
to the smallest different vertex; its intersection rank from $v$ is
$u=r$. For a quadrilateral of type $j$, choose the directed edge from
vertex zero to the smallest vertex outside $\{0,j+1\}$; the intersection
rank from zero is $u=t_0+r$. If the resulting global edge has weight
$W_e$, interval offset $O_e$, and the directed local edge agrees with
its global orientation when $\epsilon=0$, then the marker is
\[
 p=O_e+
 \begin{cases}
 u,&\epsilon=0,\\
 W_e-1-u,&\epsilon=1.
 \end{cases}
\]
The source edge-orientation data include the reversal needed if the
chosen directed edge opposes the canonical local edge ordering.

\begin{lemma}[Affine marker correctness]
\label{lem:affine-markers}
The preceding formulas address one vertex of the marked normal disc
and the two actual cut chambers adjacent to its opposite sides.
Their evaluation is polynomial in source size and coordinate bit length
and does not depend on the disc's numerical rank through iteration.
\end{lemma}
\begin{proof}
Nested triangles at one vertex cut off the consecutive arm chambers;
after the last triangle, the next chamber is the central chain endpoint
on that vertex's side of the quadrilaterals. Consecutive quadrilaterals
separate consecutive chain chambers. These observations give the two
chamber formulas, including a zero-quadrilateral tetrahedron.

Along an edge starting at $v$, the $t_v$ triangle intersections occur
first, in their vertex-stack order. Along an edge crossing a quadrilateral
type from vertex zero, the quadrilateral intersections follow the
$t_0$ triangle intersections in their chain order. Hence the indicated
rank is a vertex of the specified disc. Matching equations identify the
same globally oriented intersection ranks across every source face.
Reversal of a finite ordered set of $W_e$ intersection points sends
$u$ to $W_e-1-u$, proving the last formula. All expressions use a bounded
number of additions, comparisons, and orientation lookups after the
source offset scan, per representative. They perform no loop over $q$,
$t_v$, or $r$.
\end{proof}

The independent checker derives these side addresses through horizontal
faces of its raw cell charts; it does not copy the exporter's formulas
as an assumed list of markers.

For component $i$, let $p_i^-$ and $p_i^+$ be the two chamber atoms adjacent
to its representative disc. Assign them scalar weights $2^{2i}$ and
$2^{2i+1}$. If several side markers name the same atom, their weights add.

\begin{theorem}[Source-authenticated chamber export]
\label{thm:native-chamber}
Assume $k>0$ and that the representative census and the two-sidedness gate
both succeed.
The exact chamber-side histogram determines the oriented chamber graph of
the labelled embedded normal surface. Its vertices are the histogram masks;
edge $i$ runs from the mask containing bit $2i$ to the mask containing bit
$2i+1$. Every mask is nonzero and has multiplicity one, the masks are
pairwise bit-disjoint, and together they contain all $2k$ marker bits.
The resulting graph is connected.
\end{theorem}
\begin{proof}
For a connected two-sided component $F_i$, each of
$F_i\times\{-1\}$ and $F_i\times\{+1\}$ is connected and lies in one
chamber of the full cut. Thus a marker at any disc on that side identifies
the chamber containing the entire side.

Every chamber touches at least one side when the surface is nonempty.
Otherwise that chamber has no surface boundary along which it can be glued
to another chamber. It would remain a separate component when all cut sides
are reglued, contrary to connectedness of the source. Because every side has
one authenticated marker, no chamber is unmarked.

The orbit sum on a chamber is exactly the sum of its distinct side-marker
powers. Hence it is a nonzero mask. No two chambers can contain the same
bit, so nonzero masks cannot repeat and their multiplicities are one.
Every marker lies in a chamber, giving exhaustive bit coverage. The stated
edge construction records exactly the two chambers adjacent to each $F_i$.
Regluing all components recovers the connected source, so this graph is
connected.
\end{proof}

The native exporter also checks these consequences explicitly: the bit
partition, complete coverage, and graph connectedness. It retains loops,
multiple edges, and bridges. Vertex labels are made deterministic by sorting
the chamber masks according to their smallest present bit. Empty coordinates
give a separate exact case with one vertex, no edges, and mask zero, using
the validated connected-source hypothesis.

\begin{corollary}[All labelled subcuts from one native graph]
\label{cor:batch-optimizer}
Let $G$ be an accepted chamber graph and let $J$ be a subset of its surface
labels. Then
\[
 \#\pi_0\bigl(M\mathbin{\backslash\!\backslash}
                \textstyle\bigcup_{i\in J}F_i\bigr)
   =\#\pi_0(G\setminus J).
\]
Consequently, $J$ is simultaneously nonseparating if and only if
$G\setminus J$ is connected. The largest such batch has cardinality
$k-|V(G)|+1$, and the complement of a spanning tree attains this value.
For nonnegative additive retention weights, a maximum-reward batch is the
complement of a minimum-weight spanning tree.
\end{corollary}
\begin{proof}
Start with the full cut and glue back the components whose labels do not
belong to $J$. Gluing the two sides of one connected component merges exactly
the chamber components at the endpoints of its graph edge; if the endpoints
already agree, the component partition is unchanged. After all such gluings,
the remaining connected-component partition is precisely that of the
subgraph containing the edges outside $J$. A connected spanning subgraph has
at least $|V(G)|-1$ edges. A spanning tree reaches this minimum; minimizing
its deleted retention weight gives the weighted statement.
\end{proof}

The graph-topological and cographic consequences in this corollary are
already present at the abstract level in report~92. The new work supplies
their independently authenticated native source input on the stated domain.

\begin{figure}[htbp]
\centering
\begin{tikzpicture}[
 chamber/.style={circle,draw=ink,fill=pale,minimum size=11mm},
 every edge/.style={draw=ink,thick},
 >=Stealth]
\node[chamber] (a) at (0,1.8) {$C_0$};
\node[chamber] (b) at (-1.8,-0.8) {$C_1$};
\node[chamber] (c) at (1.8,-0.8) {$C_2$};
\draw[thick,->,ink] (a) -- node[left] {$F_0$} (b);
\draw[thick,->,ink] (b) -- node[below] {$F_1$} (c);
\draw[thick,->,accent] (c) -- node[right] {$F_2$} (a);
\node[align=left,text width=6.6cm,anchor=west] at (3.1,0.4)
 {Three parallel meridian discs produce three chambers and a cycle.\\[0.7em]
 Any one disc is a nonseparating batch: deleting one edge leaves the graph
 connected. Deleting two edges disconnects it.};
\end{tikzpicture}
\caption{A native stored example. Edge orientations choose component
coorientations and do not affect the connectivity criterion. The highlighted
edge is the complement of a spanning tree.}
\label{fig:three-meridians}
\end{figure}

\begin{proposition}[The chamber cycle obstruction, recalled]
\label{prop:cycle-obstruction}
For a connected orientable source and a disjoint system of connected
two-sided properly embedded surfaces,
\[
 \dim_{\Q}H_1(G;\Q)=k-|V(G)|+1
 \ \le\ \dim_{\Q}H_1(M;\Q).
\]
\end{proposition}
\begin{proof}
Choose a spanning tree of $G$. For every edge outside the tree, its edge
together with the tree path between its endpoints is a fundamental cycle;
loops are included. Realize the crossings of that cycle in small collars
of the corresponding surface components. Within each intervening chamber,
join consecutive collar endpoints by a path, using connectedness of the
chamber. The resulting closed curve in $M$ crosses the chosen non-tree
component once and crosses no other non-tree component.

Coorient the surfaces and orient each curve accordingly. Algebraic
intersection with the non-tree components defines linear functionals on
$H_1(M;\Q)$: the surfaces determine relative homology classes and their
boundaries lie in $\partial M$. On the constructed curves the matrix of
these functionals is the identity. Thus these curves have linearly
independent homology classes, one for each of the $k-|V(G)|+1$ non-tree
edges.
\end{proof}

This proof establishes existence of the paths. It supplies neither a
compressed path representation nor a source-support certificate for a
later edit. That distinction matters when using the graph in a hierarchy.
The obstruction itself is earlier mathematics from report~92.

\subsection{Independent verification and encoding costs}

The certificate \texttt{normal-chamber-graph-v1} contains the source
fingerprint, disc representatives, all marker addresses, the graph, and
three proof records: a weighted surface-label census, an unweighted
orientation-double count, and a weighted chamber-side census. Source changes
or changed normal coordinates invalidate the fingerprint.

The checker obtains the two chamber atoms from the independent raw cut-cell
reference~\cite{cutcells}, using its horizontal-face-to-chamber incidences. It derives the
surface marker from an independently reconstructed normal-arc recipe,
including the source edge orientation. It independently reconstructs chamber
pairings through the face-region reference, then replays all three orbit
proofs. It imports no exporter, chamber producer, graph decoder, or orbit
search. The validated finite source and the established normal-arc compiler
remain shared primitives, as in the maintained topology checkers.

Let $T$ be the tetrahedron count, $B$ the maximum normal-coordinate
bit length, and $k$ the explicit label count. The source incidence data
use $O(T\log(T+1))$ bits. Each orbit system has $O(T)$ interval
pairings with polynomially bounded endpoint bit lengths. Each marker query
starts with $O(k)$ runs. All marker weights are nonnegative and their total
mass is less than $2^{2k}$, so all transported scalars have $O(k)$ bits.
The graph has $k$ edges and at most $k+1$ vertices. Its incidence lists use
$O(k\log k)$ bits, while the current integer-mask and weight-run records
can require $O(k^2)$ bits. The complete construction and replay remain
polynomial in $T$, $B$, $k$, and the certificate trace lengths.

This is an output-sensitive assertion. A binary vector can describe
$N$ parallel discs, with $k=N$ exponential in the coordinate encoding.
Explicitly labelling all those components requires a correspondingly large
output. The operator does not hide this cost or silently quotient parallel
components. Incomplete component coverage is refused. In particular,
large-coordinate successful examples below are connected surfaces with
exponentially many discs, rather than exponentially many labelled components.

All three orbit queries share the caller's cycle allowance. An exhausted
allowance yields \texttt{INCONCLUSIVE} without a graph or certificate.
The checker can bound the combined number of recorded orbit events.
Neither limit is represented as a topological negative conclusion.

\subsection{An explicit native interface}

The supplied triangulation uses the maintained finite-source JSON schema;
normal coordinates are seven entries per tetrahedron in the established
four-triangle, three-quadrilateral order. Disc columns and ranks begin at
zero. A minimal use of the new API is:

\begin{lstlisting}[language=Python]
from fastunknot.normal_chamber_graph import normal_chamber_graph
from fastunknot.normal_chamber_graph_verify import verify_normal_chamber_graph

representatives = [{"tetrahedron": t, "column": c, "rank": r}]
result = normal_chamber_graph(
    triangulation, coordinates, representatives,
    record_certificate=True,
)
# Authenticate before consuming a successful graph:
if result["status"] == "COMPLETE":
    accepted = verify_normal_chamber_graph(
        triangulation, coordinates, result["certificate"]
    )
\end{lstlisting}

The complete call and checker signatures, including work bounds and
callback arguments, appear in the delivered source and reproduction
examples. The graph result is a new observer; it does not dispatch an
unknot verdict or replace the supplied-cut handlebody checker.

